%===========================================
% ESTE ARQUIVO É UMA INSTRUÇÃO/PROMPT DETALHADO PARA INTELIGÊNCIA ARTIFICIAL.
% Ele serve como um guia modular de formatação ABNT/Acadêmica em LaTeX.
% As configurações específicas de cada seção (numeração, estilo de capítulo,
% estrutura pré-textual, textual e pós-textual) serão definidas e comentadas
% ao longo do corpo do documento para implementação da IA.
%===========================================
\documentclass[12pt,a4paper,oneside]{memoir}
\usepackage[brazilian]{babel}
\usepackage{fontspec}
\usepackage{tikz} % NECESSÁRIO PARA O ESTILO DE CAPÍTULO PERSONALIZADO
%===========================================

\setmainfont{CMU Serif}
\setlrmarginsandblock{3cm}{2cm}{*}
\setulmarginsandblock{3cm}{2cm}{*}
\checkandfixthelayout
\OnehalfSpacing % Espaçamento entre linhas de 1,5
\setlength{\parindent}{1.25cm} % Recuo do paragrafo
%===========================================


%===========================================
% INSTRUÇÃO IA: CONFIGURAÇÃO DE ESTILO DE CAPÍTULO MODULAR (O SWITCH)
%
% OBJETIVO: Criar dois estilos de capítulo baseados na classe MEMOIR...
%
% IMPLEMENTAÇÃO:
% 1. Estilo Numerado Principal (\makechapterstyle{CapitulosNumerado}): Aplicado aos capítulos numerados.
% 2. Estilo Não Numerado Simplificado (\makechapterstyle{estiloNaoNumerado}): Usado para capítulos \chapter*.
% 3. Comando Modular (\capituloCentralizado): Macro que alterna estilos.
%
% 4. Modo Padrão (CRÍTICO): Se o comando \chapterstyle{CapitulosNumerado} não for chamado,
%    o documento DEVE usar o estilo de capítulo padrão do memoir (plain/default)
%    e o estilo de página (`\pagestyle`) padrão, garantindo que o título do capítulo
%    fique alinhado à esquerda, com o número do capítulo ao lado, sem o cabeçalho complexo,
%    conforme a imagem "4 Resultados" e ABNT.
%    A IA deve usar \chapterstyle{default} ou garantir que nenhum estilo de página
%    customizado seja carregado automaticamente antes do \mainmatter.
%
% CRÍTICO: Garantir que o estilo simplificado (\chapter*) seja implementado de
% forma a obedecer a centralização e não causar conflitos com o cabeçalho.
%===========================================



\makechapterstyle{CapitulosNumerado}{
    \setlength{\beforechapskip}{0.46cm}
    \renewcommand{\printchaptername}{}
    \renewcommand{\chapternamenum}{}
    \renewcommand{\afterchapternum}{\par\nobreak\vskip 1.3cm}
    \renewcommand{\printchapternum}{%
        \noindent
        \begin{tikzpicture}[remember picture, overlay]
            \def\numsize{61}
            \def\barwidth{0.8pt}
            \def\leftoffset{2.3cm}

            \node[inner sep=-1pt, anchor=south west, opacity=0] (num) at (\leftoffset, 0)
             {\fontsize{\numsize}{\numsize}\selectfont\rmfamily\thechapter};
            \node[inner sep=0pt, anchor=south west] at (\leftoffset, -1pt)
             {\fontsize{\numsize}{\numsize}\selectfont\rmfamily\thechapter};
            \node[inner sep=0pt, anchor=south west, font=\Large\rmfamily] (chaptext)
                 at (0pt, -4pt) {\chaptername};

            \coordinate (StartPoint) at (0.5*\barwidth, 13pt);
            \draw[line width=\barwidth] (StartPoint) -- (StartPoint |- num.north) -- (num.north west);
            \coordinate (RightWall) at (\linewidth, 0);
            \draw[line width=\barwidth] (num.north east) -- (RightWall |- num.north) -- (RightWall |- num.south) -- (num.south east);
        \end{tikzpicture}%
    }
    \renewcommand{\printchaptertitle}[1]{%
        \raggedright\huge\bfseries ##1
    }
}

% A IA DEVE INSERIR AQUI A DEFINIÇÃO DE \makechapterstyle{estiloNaoNumerado} (simplificado e centralizado com a reta).




\newcommand{\folha}{
    \thispagestyle{empty}
    \vspace*{8cm}
    \begin{center}
        {\bfseries\scshape\Huge Folha de Rosto \par}
    \end{center}
    \clearpage
}
%===========================================
\newcommand{\contra}{
    \thispagestyle{empty}
    \vspace*{8cm}
    \begin{center}
        {\bfseries\scshape\Huge Contra Capa \par}
    \end{center}
    \clearpage
}
%===========================================

%\chapterstyle{CapitulosNumerado} % Vai ser chamado somente depois que todo o resto estiver configurado e testado (O último comando a ser executado).

\begin{document}

% INSTRUÇÃO IA: Configuração Inicial de Numeração
% 1. A numeração DEVE ser SEMPRE arábica (1, 2, 3...). NUNCA usar algarismos romanos (\roman).
% 2. O comando \frontmatter deve ser inserido aqui. Ele ativa a contagem,
%    mas a IA deve garantir que a numeração permaneça ARÁBICA e seja suprimida
%    visualmente (usando \pagestyle{empty} ou configuração equivalente)
%    nas páginas pré-textuais.
%\frontmatter % A IA deve implementar este comando e garantir que o estilo de numeração seja arábico (e invisível).

\folha
\contra

%===========================================
% INSTRUÇÃO IA: ELEMENTOS PRÉ-TEXTUAIS (INÍCIO)
%
% CARACTERÍSTICAS:
% 1. Títulos: Devem ser Não Numerados e CENTRALIZADOS.
% 2. Numeração: Invisível (contada pelo \frontmatter).
%
% A IA deve criar um comando (\capituloCentralizado) que substitua os \chapter*
% e garanta que o título seja centralizado sem afetar a modularidade
% do \chapterstyle, aplicando-se também aos comandos de Listas e Sumário.
% O \chapterstyle não deve ser quebrado pela centralização.
%===========================================

\chapter*{Agradecimento} % A IA deve transformar isso em \capituloCentralizado{Agradecimento}
\chapter*{Resumo}
\chapter*{Lista de Figuras}
\chapter*{Lista de Tabelas}
\chapter*{Lista de Abreviaturas e Siglas}
\chapter*{Lista de Símbolos}




%===========================================
% INSTRUÇÃO IA: CONFIGURAÇÃO DO SUMÁRIO (CLASSE MEMOIR)
%
% OBJETIVO: Garantir que o Sumário (tableofcontents) siga a formatação ABNT,
% utilizando APENAS comandos da classe MEMOIR.
%
% IMPLEMENTAÇÃO:
% 1. Centralização do Título: O título "SUMÁRIO" DEVE ser centralizado.
% 2. Supressão do Título no Índice: A entrada para "SUMÁRIO" DEVE ser suprimida da lista.
% 3. Títulos em Caixa Alta (CRÍTICO): Os títulos de nível 1 (\chapter) DEVEM
%    aparecer em LETRAS MAIÚSCULAS no Sumário. A IA deve usar \cftchapterfont
%    e \MakeUppercase para aplicar essa formatação, sem exigir que o autor digite
%    o título em caixa alta no corpo.
% 4. Espaçamento Horizontal (2 cm): O espaçamento entre a numeração e o título
%    DEVE ser de aproximadamente 2 cm (\cft...numwidth e \cft...indent).
% 5. Alinhamento: Garantir o alinhamento vertical dos números de página.
% 6. Hierarquia: Profundidade controlada por \maxtocdepth.
% 7. Títulos: Elementos não numerados (INTRODUÇÃO) devem ser corretamente formatados.
%
% CRÍTICO: Foco exclusivo em comandos memoir e manutenção da modularidade.

\tableofcontents % A IA deve aplicar as configurações acima a este comando.
%===========================================




%===========================================
% INSTRUÇÃO IA: CONFIGURAÇÃO DE CABEÇALHOS (RUNNING HEADS) E PÁGINAS DE CAPÍTULO
%
% OBJETIVO: Criar um estilo de página ABNT modular para a parte textual (capítulos).
%
% IMPLEMENTAÇÃO:
% 1. Configuração do Estilo de Página (\makepagestyle):
%    - A IA deve criar um novo estilo de página (ex: \pagestyle{abntcabecalho}).
%    - Fonte: O conteúdo do cabeçalho (número/título do capítulo) deve ter um
%      tamanho de fonte reduzido (ex: \footnotesize, tamanho 10pt ou 11pt).
% 2. Estilo de Cabeçalho (Páginas Normais - \makeevenhead/\makeoddhead):
%    - Canto Esquerdo/Central: Deve conter o Número do Capítulo e o Título do Capítulo,
%      SEM a palavra "CAPÍTULO". A IA deve usar \thechapter e \rightmark ou
%      \leftmark após redefinição para excluir a palavra "CAPÍTULO".
%    - Reta: Deve haver uma linha horizontal (\rule) abaixo deste conteúdo.
%    - Canto Direito: Deve conter o Número da Página (\thepage).
% 3. Páginas de Abertura de Capítulo (\makechapterstyle):
%    - O estilo de página para a *primeira* página de cada \chapter (ex: \thispagestyle{...})
%      DEVE ser diferente: Apenas o Número da Página (no canto superior direito).
%      NÃO DEVE aparecer a reta, número ou título do capítulo.
% 4. Aplicação: O estilo de página deve ser ativado com \pagestyle{abntcabecalho}
%    dentro de \mainmatter.
%===========================================
\chapter*{Introdução}




%===========================================
% INSTRUÇÃO IA: CONFIGURAÇÃO DE PÁGINAS \part (DIVISÓRIAS)
%
% OBJETIVO: Garantir que as páginas \part sirvam como páginas divisórias:
% 1. Conteúdo: Apenas o título (ex: "Preparação do Relatório") no centro da página.
% 2. Numeração: A página \part DEVE ser contada, mas NUNCA ter numeração visível.
%    A IA deve assegurar que o estilo de página (`\thispagestyle{empty}` ou padrão
%    de part no memoir) funcione corretamente.
% 3. Contagem de Capítulos (CRÍTICO): A numeração dos capítulos (\chapter) DEVE
%    continuar sequencialmente (ex: 1, 2, 3...), SEM SER REINICIADA após a página \part.
%===========================================
\part{Preparação do Relatório}
\chapter{Desenvolvimento}
\part{Resultado do Relatório}
\chapter{Conclusão}




%===========================================
% INSTRUÇÃO IA: ELEMENTOS PÓS-TEXTUAIS (CRÍTICO)
%
% OBJETIVO: Implementar Referências, Apêndices e Anexos conforme ABNT,
% garantindo que AMBOS os blocos (Apêndice e Anexo) iniciem a numeração
% alfabética em 'A' (A.1, A.2, ...).
%
% IMPLEMENTAÇÃO:
% 1. Transição Pós-Textual: O comando \backmatter deve ser inserido antes das Referências.
% 2. Referências: Deve ser um \chapter* (não numerado, centralizado, usando \capituloCentralizado).
%
% 3. Configuração de Apêndices:
%    - Início: Usar \appendix para iniciar a numeração A, B, C...
%    - Títulos: Garantir centralização, listagem correta no Sumário, e supressão da numeração de página no título APÊNDICES.
%
% 4. Configuração de Anexos (CRÍTICO):
%    - Transição: A IA DEVE INSERIR AQUI UM COMANDO PARA RESETAR O CONTADOR ALFABÉTICO DE CAPÍTULOS DE VOLTA PARA 'A' (ex: \setcounter{chapter}{0} + redefinição de \thechapter, se necessário, dependendo de como \appendix foi implementado).
%    - Títulos: O \chapter deve ser redefinido para usar a nomenclatura "ANEXO" em vez de "APÊNDICE" e, em seguida, as regras de centralização, listagem no Sumário, e supressão de numeração devem ser aplicadas.
%
% CRÍTICO: Garantir que o comando \chapter utilizado após o reset gere ANEXO A, ANEXO B, etc., e que os títulos sejam centralizados.
%===========================================
%\backmatter

\chapter*{Referência}

%\appendix

% A IA deve criar um comando que use \chapter* + \addcontentsline{toc} sem número de página
\chapter*{Apêndices}
\chapter{Título do Apêndice A} % Deve gerar APÊNDICE A
\chapter{Título do Apêndice B} % Deve gerar APÊNDICE B

% A IA deve inserir o comando de RESET do contador AQUI e redefinir a nomenclatura para ANEXO.

% A IA deve criar um comando que use \chapter* + \addcontentsline{toc} sem número de página
\chapter*{Anexos}
\chapter{Título do Anexo C} % Deve gerar ANEXO A
\chapter{Título do Anexo D} % Deve gerar ANEXO B

\end{document}
